\subsection{与三篇参考文献的对照及后续改进}
da Silva等\cite{silva2014}采用圆柱有限体积模型，将水分损失与表面汽化潜热耦合，并比较不同热方程简化。在其58.6~$^\circ$C香蕉试验中，忽略潜热使温度拟合的 $R^2$ 从0.98843降至0.75502。这是特定材料的实验结果，不能直接换算成本题误差，却表明不应仅以题面未给潜热为由认定蒸发冷却可以忽略。

后续可在完成界面含湿关系与干物质质量标定后，比较无潜热和有潜热两个模型。以向外水质量通量 $j_w>0$ 表示蒸发，则表面能量平衡为
\begin{equation}
-k\partial_nT=h_T(T_s-T_e)+L_v(T_s)j_w,\qquad
j_w=\rho_d h_m(C_s-C_{\mathrm{eq}}).\label{eq:latent_extension}
\end{equation}
式中 $L_v$ 可先使用经温标核验的自由水近似，但 $C_{\mathrm{eq}}$ 应来自材料平衡吸湿关系。还需连同水分携带的焓和经验储热关系一起检查，不能仅将 $\rho c_p\partial_tT$ 机械替换为 $\partial_t(\rho c_pT)$。本次主结果未引入该扩展，避免把尚未标定的模型当作已验证答案。

张继凯等\cite{zhang2024}对6~mm山药片的红外联合热风干燥进行了收缩与非收缩对照，并以实测温度、含水率和品质指标验证。这支持保留本题的同物性几何对照；但本题没有红外热源、局部面积修正系数和品质测量，因此不直接加入其辐射项、不重复叠加其经验面积修正，也不采用其60~$^\circ$C品质结论作为本题最佳温度。

Tzempelikos等的会议稿\cite{quince2014}使用厚度方向一维模型和外部流场计算得到的迎风、背风传递系数，含水率方程显式区分干固体密度与体积密度。该文提示应辨明几何主导方向、质量基准和外部边界；其Halsey吸湿关系及quince物性不移植到本题药材。若补充环境相对湿度、风速及内部实验数据，可再标定界面关系并检验恒定传递系数假设。

三篇文献提供了扩展模型和实物验证的依据；本文空间加密、守恒残差和二维对照属于数值验证，不能等同于这些文献的实验拟合证据。
